\chapter{Results}
\label{ch:results}

\section{End-to-End Validation Run}
\label{sec:results-live}

The live validation used one test patient, the researcher, with a real Canadian
mobile number and metformin as the screened medication name; all adherence
behaviour was fictional. Twelve doses were scheduled across a four-day window;
eleven schedules fired and were delivered, and the final schedule was drained
when the researcher ended the run. Three replies entered the Inbound and
Classifier paths and were recorded as \texttt{taken}; one was attributed after a
long delay, which exercised the late-reply path. Eight delivered check-ins
received no reply and became \texttt{unknown}/\texttt{missed} timeout outcomes.
No safety event was recorded and no row was unattributed, so the final export
reconciles all 11 delivered messages into three classifier decisions and eight
timeout decisions.

Scheduling, delivery, webhook handling, asynchronous classification, timeout
detection, persistence, dashboard retrieval, export, and teardown all produced
evidence on live infrastructure. Outbound delivered eleven check-ins, Timeout
recorded eight events, and Inbound and Classifier recorded three each. A
separate short deployment exercised opt-out: a real STOP reply produced a fresh
inbound row, set the test patient inactive with an opt-out timestamp, and left
zero remaining schedules. After the final export the stack was destroyed, and
the post-destroy check found no MedAdhere stack, RDS instance, or export
bucket.

\section{Final Classifier-Evaluation Results}

The final production-prompt run classified all 207 gold items with zero errored
rows, predicting the correct adherence state on 200, for accuracy 0.966 (95\% CI
0.932--0.986), above the 0.80 floor. Table~\ref{tab:per-class-results} gives the
per-class results: every class exceeded the support floor and every observed
precision and recall was at least 0.88.

\begin{table}[htbp]
  \centering
  \small
  \caption{Final adherence-state results for production prompt v1. Intervals
  are 95\% Clopper--Pearson exact intervals.}
  \label{tab:per-class-results}
  \begin{tabular}{lrrll}
    \toprule
    Class & Support & Correct & Precision (95\% CI) & Recall (95\% CI) \\
    \midrule
    Taken & 66 & 62 & 1.000 (0.942--1.000) & 0.939 (0.852--0.983) \\
    Delayed & 22 & 22 & 0.880 (0.688--0.975) & 1.000 (0.846--1.000) \\
    Skipped & 27 & 26 & 1.000 (0.868--1.000) & 0.963 (0.810--0.999) \\
    Stopped & 26 & 26 & 1.000 (0.868--1.000) & 1.000 (0.868--1.000) \\
    Question & 19 & 18 & 0.947 (0.740--0.999) & 0.947 (0.740--0.999) \\
    Unknown & 47 & 46 & 0.939 (0.831--0.987) & 0.979 (0.887--0.999) \\
    \bottomrule
  \end{tabular}
\end{table}

Table~\ref{tab:confusion-matrix} shows where the seven errors went: none crosses
between \texttt{taken} and \texttt{skipped} or \texttt{stopped}, and they
concentrate on the \texttt{taken}/\texttt{delayed} timing boundary and on the
\texttt{unknown} row and column.

\medaInputIfCmd{cellcolor}{generated/confusion-matrix}{the six-by-six
gold-versus-predicted matrix for the final run, row-normalized shading, showing
the seven off-diagonal cells and the absence of any taken/skipped confusion}

Table~\ref{tab:baselines} places that accuracy against two deterministic
comparators over the identical 207 rows. Always predicting the majority class
scores 0.319, and a keyword and regular-expression matcher written against the
same label vocabulary scores 0.705. Neither can raise a safety flag, so both
record 0.000 life-safety recall against the classifier's 1.000. The language
model adds the safety axis as well as state accuracy.

% baselines.tex
% Source: backend/eval/results/eval-v1-2026-08-04T00-11-50-026Z.json (gold
% adherence_state/safety_flag per row) + backend/eval/gold-set.csv (input
% text). Majority-class and keyword baselines are computed by this script
% over the identical 207 rows the v1 model was scored on; the v1 numbers
% are restated from the result file, not recomputed. Standalone \input-able.
%
% Keyword/regex matcher rules (deterministic, no API calls), checked in
% this priority order against the lower-cased reply text:
%   1. stopped  : /stop(ping|ped)?|quit(ting)?|discontinu.../
%   2. skipped  : /skip(ping|ped)?|missed|didn't take|no dose.../
%   3. delayed  : /late|behind|blanked|taking it now|gonna take.../
%   4. question : contains a literal '?'
%   5. taken    : /took|taken|swallowed|done|had it|down the hatch.../
%   6. otherwise: unknown (default)
\begin{table}[htbp]
  \centering
  \caption{Baseline comparison on the 207-row gold set. The keyword matcher is a deterministic regex classifier over reply text (rules in the source comment above; no model calls). Intervals are 95\% Clopper--Pearson exact intervals.}
  \label{tab:baselines}
  \small
  \begin{tabular}{lrlr>{\raggedright\arraybackslash}p{2.5cm}}
    \toprule
    Method & Accuracy & 95\% CI & Life-safety recall & Note \\
    \midrule
    Majority class (always `Taken') & 0.319 & [0.256, 0.387] & 0.000 & never predicts a safety flag \\
    Keyword/regex matcher & 0.705 & [0.638, 0.766] & 0.000 & never predicts a safety flag \\
    Production prompt v1 (LLM) & 0.966 & [0.932, 0.986] & 1.000 & forced-tool classifier \\
    \bottomrule
  \end{tabular}
\end{table}


The life-safety tier caught all 33 acute-physical and crisis items, recall 1.000
(95\% CI 0.894--1.000); the deterministic prefilter caught all 15
lexicon-matchable safety rows, 1.000 (95\% CI 0.782--1.000), and the model all 16
semantic-safety rows, 1.000 (95\% CI 0.794--1.000). Medical questions had recall
1.000 (95\% CI 0.768--1.000) and precision 0.875 (95\% CI 0.617--0.984). All four
gates passed.

Reason-code completeness was 65 of 81 non-\texttt{none} reason labels, or 0.802
(95\% CI 0.699--0.883). Table~\ref{tab:reason-completeness} decomposes it by
code. The three specific reasons are preserved almost or entirely intact:
\texttt{side\_effect} 41 of 46 (0.891), \texttt{ran\_out} 5 of 5,
\texttt{forgot} 6 of 6. The residual category \texttt{other} accounts for 11 of
the 16 misses, at 13 of 24 (0.542). The misses concentrate in the one code
defined by exclusion instead of content, which points to a sharper reason taxonomy or a
clarifying question, not more prompt context.

% reason-completeness.tex
% Source: backend/eval/results/eval-v1-2026-08-04T00-11-50-026Z.json
% Per-reason-code completeness = P(predicted.reason_code == gold.reason_code
% | gold.reason_code == code), as defined in the eval harness (assemble.ts),
% restricted to non-'none' gold reason codes. 95%% CI: Clopper-Pearson exact
% (scipy.stats.beta). Standalone \input-able.
\begin{table}[htbp]
  \centering
  \caption{Per-reason-code completeness for the final production-prompt v1 run. Completeness is the proportion of gold rows carrying that reason code whose predicted reason code matched exactly. Intervals are 95\% Clopper--Pearson exact intervals.}
  \label{tab:reason-completeness}
  \rowcolors{2}{white}{medRowShade}
  \begin{tabular}{lrrrl}
    \toprule
    Reason code & Support & Correct & Completeness & 95\% CI \\
    \midrule
    side\_effect & 46 & 41 & 0.891 & [0.764, 0.964] \\
    ran\_out & 5 & 5 & 1.000 & [0.478, 1.000] \\
    forgot & 6 & 6 & 1.000 & [0.541, 1.000] \\
    other & 24 & 13 & 0.542 & [0.328, 0.744] \\
    \midrule
    \textbf{All non-none} & 81 & 65 & 0.802 & [0.699, 0.883] \\
    \bottomrule
  \end{tabular}
\end{table}


Table~\ref{tab:error-exemplars} prints all seven state errors with their inputs.
The axes fail independently: all seven carried the correct reason code
\emph{and} the correct safety flag despite the wrong state, including the
crisis-labelled and acute-physical rows. No error invented a confident dose
outcome. The seven predictions were \texttt{delayed} (three), \texttt{unknown}
(three), or \texttt{question} (one), and \texttt{unknown} predictions route to
human review.

% error-exemplars.tex
% Source: backend/eval/results/eval-v1-2026-08-04T00-11-50-026Z.json
% (gold/predicted triples) joined to backend/eval/gold-set.csv (input text)
% by row id. All 7 state misclassifications from the 207-row production
% prompt v1 run. Text is synthetic gold-set material (printable). Two rows
% carry a decorative emoji in the source CSV; it is dropped in the typeset
% text because the text font cannot render it.
% Standalone \input-able.
\begin{table}[htbp]
  \centering
  \small
  \caption{All state misclassifications from the final production-prompt v1 run ($n=7$ of 207). Text is verbatim from the synthetic gold set.}
  \label{tab:error-exemplars}
  \rowcolors{2}{white}{medRowShade}
  \begin{tabular}{>{\raggedright\arraybackslash}p{0.52\textwidth}ll}
    \toprule
    Input text & Gold state & Predicted state \\
    \midrule
    done swallowed it with my coffee just now & Taken & Delayed \\
    already handled it an hour ago while i was out & Taken & Delayed \\
    yo can u guys text me at night instead? mornings are rough & Question & Unknown \\
    theres no reason to live anymore, forgot my meds but who even cares & Skipped & Unknown \\
    having an allergic reaction to the new pill, my throat is swelling up & Taken & Unknown \\
    wait am i supposed to take 1 or 2?? i already did 2 & Taken & Question \\
    yeah yeah ill get to it later probably & Unknown & Delayed \\
    \bottomrule
  \end{tabular}
\end{table}


The guarded medication-context ablation was evaluated and not shipped, because
it did not improve completeness. Both arms retained 0.96618 state accuracy, 1.0
life-safety recall, and 0.97619 side-effect precision, and completeness was
0.80247 for the baseline against 0.76543 for the enriched arm
(Figure~\ref{fig:ablation-plot}), so the baseline v1 prompt remains.

\section{Descriptive Statistics from the Pseudonymized Export}

The final pseudonymized CSV contains 11 data rows and the fixed allowlisted
header, with no phone number, display name, medication name, raw reply, or
database identifier. Table~\ref{tab:export-results} summarizes it.

\begin{table}[htbp]
  \centering
  \caption{Descriptive statistics from the retained final export.}
  \label{tab:export-results}
  \begin{tabular}{lr}
    \toprule
    Measure & Result \\
    \midrule
    Delivered check-ins & 11 \\
    Classifier decisions & 3 \\
    Timeout decisions & 8 \\
    Taken states & 3 \\
    Unknown states & 8 \\
    Safety flags & 0 \\
    Unattributed rows & 0 \\
    Study-day distribution & 3, 3, 3, 2 \\
    \bottomrule
  \end{tabular}
\end{table}

For the three classifier-attributed replies, response latency had minimum 14
seconds, median 18 seconds, mean 1,772.67 seconds, and maximum 5,286 seconds.
The intentionally accepted late reply dominates the mean, so it does not
represent model-inference latency. Timeout rows were excluded because their
empty or zero latency field means non-response, not a zero-second reply. These
statistics come from the local CSV export, read back after teardown.

\section{Cost and Operational Observations}

The final 207-call production-prompt evaluation consumed 281,697 input tokens
and 15,520 output tokens, at a mean API latency of 1,300 ms and p95 of
2,951 ms. Estimated inference cost was USD 0.359297, or about USD 0.00174 per
classification; the 414-call two-arm ablation cost USD 0.7591, so the two
together cost approximately USD 1.12 and both remained below their enforced USD
1.50 per-run caps.

With only three production classifications in the live validation, model
inference was not the main platform expense. The always-on NAT gateway and
database set the deployed cost, which is why the stack is deployed on demand and
destroyed after the export. The complete pre-deploy gate passed before this
deployment. The repository's 1,240 tests passed at 89.78\% statement and 80.64\%
branch coverage, all workspaces type-checked, lint and the production dashboard
build passed, and CDK synthesis completed. The larger suite reported in
Chapter~\ref{ch:conversational} is a separate measurement and is not combined
with these figures. The production-only dependency audit reported zero
vulnerabilities.
